\documentclass[compress]{beamer}    %slidestop=title in up-left corner (with "slidescentered" centered); compress=navigation bars as small as possible

\usepackage[ugm]{mathdesign} %use Garamond font

%gets rid of bottom navigation bars and puts slide numbers right bottom
\setbeamertemplate{footline}[frame number]{}

%alternative footer: name, institute, page number
%\setbeamertemplate{footline}
%{%
%  \leavevmode%
%  \hbox{\begin{beamercolorbox}[wd=.5\paperwidth,ht=2.5ex,dp=1.125ex,leftskip=.3cm plus1fill,rightskip=.3cm]{author in head/foot}%
%    \usebeamerfont{author in head/foot}\insertshortauthor
%  \end{beamercolorbox}%
%  \begin{beamercolorbox}[wd=.5\paperwidth,ht=2.5ex,dp=1.125ex,leftskip=.3cm,rightskip=.3cm plus1fil]{title in head/foot}%
%    \usebeamerfont{title in head/foot}\insertshortinstitute \hfill \insertframenumber \,/\,\inserttotalframenumber
%  \end{beamercolorbox}}%
%  \vskip0pt%
%}

%gets rid of navigation symbols
\setbeamertemplate{navigation symbols}{}


%%complete themes (not recommended)
%%\usetheme{Madrid}             %"Berkeley":lhs navigation bar; 1st alternative: "Madrid" without a navigation bar but with a footer including name (institution), title, date page numbering (but then \institute{Univ. ...}); 2nd alternative: defaults without all the extras; 3rd alternative "Dresden": two bottom lines)
%%"Dresden is possibly easy to adapt to UvT colors
%instead of using a complete theme inner and outer themes can be specified in isolation
%inner theme:
\useinnertheme{rounded} %actually not necessary but nice for theorem boxes etc. but enumerate listings look a bit stupid! Alternative: default = just comment it out
%outer themes:
%\useoutertheme[subsection=false,footline=authorinstitutetitlepage]{miniframes} %other options for footer see beamer user guide p. 153


%%complete color themes: (ok but inner outer styles below are better)
%%\usecolortheme{albatross}    % background is colored (in blue) in albatross; the option [overlystylish] installs a background canva which is definitely too stylish for an economics department
%%\usecolortheme{seagull}   %for greyscale (black and white)
%%default is not bad!!!

%alternative to complete color schemes one can define an inner and outer color scheme
%inner color themes
%\usecolortheme{orchid} %deep color in box(theorem etc.) background fits with outer color scheme "whale"
\usecolortheme{rose} % less aggressive slight grey shading background (in boxes/theorems etc.) fits nicely with outer color scheme \usecolortheme{seahorse}
%outer color theme
\usecolortheme{dolphin}   %alternative to seahorse: "whale" but only if inner scheme is "orchid"
%
%\usefonttheme{structurebold}       %titles, headlines, foodlines and sidebars are bold

\usepackage{amsmath, eurosym, textcomp, graphicx, amsthm, amssymb,amsfonts, natbib, url, xcolor, tikz, subfig}

\usepackage[english]{babel} %English language if Koma script classes are used (Bibliography instead of Inhaltsverzeichnis..)
\usepackage[latin1]{inputenc}
\usepackage[T1]{fontenc} %Encoding 
\usepackage{ae,aecompl}%avoids the problem of faint or hardly printable output


\usepackage{sgamevar} %Osborne's packages to draw games
\usepackage{hyperref}
%\usepackage[all]{hypcap}


% \usepackage{comment}
% \newcommand{\com}{\comment}
% \newcommand{\ecom}{\endcomment}

%\usepackage{ngerman}
\setlength {\parindent}{0cm}
\newcommand {\bi}{\begin{itemize}}
\newcommand {\ei}{\end{itemize}}
\newcommand{\be}{\begin{enumerate}}
\newcommand{\ee}{\end{enumerate}}
\newcommand{\Ra}{\Rightarrow}
\newcommand{\ra}{\rightarrow}
\newcommand{\Lra}{\Leftrightarrow}
\newcommand{\del}{\partial}
\newcommand{\bea}{\begin{eqnarray*}}
\newcommand{\eea}{\end{eqnarray*}}
\newcommand{\less}{<}
\newcommand{\great}{>}
\newcommand{\vbar}{|}
\renewcommand{\Re}{\mathbb{R}}

\theoremstyle{plain}
\newtheorem{prop}{Proposition}
%\newtheorem{theorem}{Theorem}
\newtheorem{assumption}{Assumption}
%\newtheorem{lemma}{Lemma}
\newtheorem{proposition}{Proposition}
\newtheorem{observation}{Observation}
%\newtheorem{corollary}{Corollary}
\theoremstyle{definition}
\newtheorem{ex}{Example}
\newtheorem{exercise}{Exercise}
%\newtheorem{definition}{Definition}

%\setlength{\parskip}{1.5ex plus 0.5ex minus 0.5ex}
\author{Christoph Schottm\"uller}
\institute{University of Copenhagen}
\title{Game Theory\\ Refinements: Perfect equilibrium}
\date{}

\newcommand{\fr}{\frame[allowframebreaks]}
\newcommand{\ft}{\frametitle}

\newcommand{\backupbegin}{
   \newcounter{framenumberappendix}
   \setcounter{framenumberappendix}{\value{framenumber}}
}
\newcommand{\backupend}{
   \addtocounter{framenumberappendix}{-\value{framenumber}}
   \addtocounter{framenumber}{\value{framenumberappendix}}
}

\begin{document}
\maketitle
\section[Outline]{}
\frame{\ft{Outline}\tableofcontents}


\section{Perfect equilibrium in strategic form games}

\fr{\ft{Refinements}
\begin{itemize}
\item multiplicity of  NE 
\item some NE ``feel a bit unreasonable''
\item ``refinement'': introduce some criterium that some NE satisfy and others do not
\item most useful in extensive form games $\Ra$ we will return to refinement later
\end{itemize}


\begin{ex}[]\label{ex:unreason_ne}
   \begin{table}[h]
\centering
%\begin{tabular}{l|c|c}
 \begin{game}{2}{2}
      \> L    \>  R\\% \hline
T   \>1,1    \> 0,0  \\
B   \>0,0    \>0,0  
\end{game}
%\end{tabular}
\caption{$(B,R)$ is NE; intuitive outcome?}\label{tab:ne_in_weak_dom_strat}
\end{table}
% The game in table \ref{tab:ne_in_weak_dom_strat} has two NE: $(T,L)$ and $(B,R)$. Somehow the $(B,R)$ equilibrium does not really feel right. For example, both players use a weakly dominated action in this equilibrium. A first refinement could be that we require equilibria do not use dominated actions (or at least one player uses an undominated action).
\end{ex}
}


\fr{\ft{Perfect equilibrium}
  \begin{itemize}
  \item problem with equilibrium $(B,R)$: other player might
    make a mistake
\item say, with some small probability $\eta$ a player ``trembles''
 \item then strictly better for player 1 (2) to choose $T$ ($L$).

\framebreak

\item formally:
  \begin{itemize}
  \item $G=\langle N,(A_i),(u_i)\rangle$ 
   \item $m_i$: number of actions of player $i$
   \item ``tremble'': vector $\eta\in
    \Re^{m_1+\dots+m_n}_{++}$ where $\eta_i^k$ is probability that
    action $a_i^k\in A_i$ is chosen accidentally by player $i$
\end{itemize}
\begin{itemize}
\item define $G'(\eta)$ as game $G$ but with strategy space such that
$$S_i(\eta)=\{\alpha_i\in \Delta A_i: \alpha_i^k\geq \eta_i^k\}$$
i.e. actions have to be \textbf{completely mixed} and
  put minimum probability $\eta_i^k$ on action $a_i^k$
\item a mixed NE exists in $G'(\eta)$ \\(proof similar to the Nash theorem)
\end{itemize}
\end{itemize}
\begin{definition}[perfect equilibrium]
  A Nash equilibrium in $G$ is (trembling hand) perfect if
\begin{itemize}
\item  there is a sequence $(\eta_t)_{t=1}^{\infty}$ of trembles converging to zero
\item  and a sequence of Nash equilibria $\alpha(\eta_t)$ in $G'(\eta_t)$ for each $t$,
\end{itemize}
 such that the sequence $(\alpha(\eta_t))_{t=1}^{\infty}$ converges to the Nash equilibrium in $G$.
\end{definition}



\framebreak

\begin{itemize}
\item previous example:
   \begin{table}[h]
\centering
%\begin{tabular}{l|c|c}
 \begin{game}{2}{2}
      \> L    \>  R\\% \hline
T   \>1,1    \> 0,0  \\
B   \>0,0    \>0,0  
\end{game}
%\end{tabular}
%\caption{Is $(B,R)$ an intuitive NE?}\label{tab:ne_in_weak_dom_strat}
\end{table}
  \begin{itemize}
  \item unique equilibrium of any $G'(\eta)$ is $(T,L)$
  \item only $(T,L)$ is perfect equilibrium
  \end{itemize}
\end{itemize}

\framebreak

\begin{theorem}
  Every finite strategic game has a perfect equilibrium.
\end{theorem}
\textbf{Proof. } (sketch)% Take an arbitrary sequence of trembles converging to the zero vector. As argued above, Nash equilibria exist in all $G'(\eta_t)$. Note that all the mixed strategy vectors $\alpha(\eta_t)$ are taken from a compact set. Hence, the sequence of Nash equilibria has a convergent subsequence. Its limit is a Nash equilibrium in the original game $G$ (this is shown as in the proof of the Nash theorem). Therefore, this limit is a perfect equilibrium.
\qed

\framebreak

\begin{ex}
  \begin{table}[h]
\centering
%\begin{tabular}{l|c|c}
 \begin{game}{2}{2}
      \> L    \>  R\\% \hline
T   \>1,1    \> 0,0  \\
B   \>1,0    \>2,1  
\end{game}
%\end{tabular}
%\caption{Is $(B,R)$ an intuitive NE?}\label{tab:ne_in_weak_dom_strat}
\end{table}
\begin{itemize}
\item Show that (T,L) is not a perfect equilibrium and (B,R) is a perfect equilibrium
\end{itemize}
\end{ex}

\framebreak

\begin{definition}[$\varepsilon $-perfect equilibrium]
  A mixed strategy profile $\alpha^\varepsilon $ is an \emph{$\varepsilon $-perfect equilibrium} if it is completely mixed and $U_i(\alpha_{-i}^\varepsilon ,a_i)\great U_i(\alpha_{-i}^\varepsilon ,a_i')$ for any $a_i,a_i'\in A_i$ implies $\alpha_i^\varepsilon (a_i')\leq \varepsilon $.
\end{definition}

\begin{itemize}
\item in short: completely mixed but only best responses are played with more than $\varepsilon $ probability
\end{itemize}

\framebreak

\begin{theorem}\label{thm:def_perf_eq}
  Let $\alpha$ be a mixed strategy profile in $G$. The following statements are equivalent
  \begin{enumerate}
  \item $\alpha$ is a perfect equilibrium
  \item $\alpha$ is a limit of $\varepsilon $-perfect equilibria for $\varepsilon $ going to zero
  \item $\alpha$ is the limit of a sequence $(\alpha(\varepsilon ))_{\varepsilon \ra 0}$ of completely mixed strategies such that $\alpha$ is a best response to $\alpha(\varepsilon )$ for $\varepsilon $ small enough.
  \end{enumerate}
\end{theorem}
\textbf{Proof. }
\begin{itemize}
\item (1) implies (2):
  \begin{itemize}
  \item If $\alpha$ is a perfect NE, there is a sequence
    $\alpha(\eta_t)$ of NE in $G'(\eta_t)$ 
    converging to $\alpha$. 
  \item\dots%  For each $\eta_t$, let $\varepsilon
  %   (\eta_t)=max_{k\in\{1,\dots,m_i\}, i\in N}(\eta_t)_i^k$.
  % \item  $\alpha(\eta_t)$ is completely mixed
  % \item As $\alpha(\eta_t)$ is a NE
  %   in $G'(\eta_t)$, any action $a_i^k$ that is not a best response
  %   will be played only with probabiltiy $(\eta_t)_i^k\leq \varepsilon
  %   (\eta_t)$
  % \item $\alpha(\eta_t)$ is an $\epsilon(\eta_t)$ equilibrium
  \end{itemize}
\item  (2) implies (3):
  \begin{itemize}
  \item let $a_i^k$ be played with positive probability under
    $\alpha$
  \item \dots % Since $\alpha$ is a limit of $\varepsilon $-perfect
  %   equilibria, $a_i^k$ has to be played with probability higher than
  %   $\varepsilon $ for $\varepsilon$ small enough
  % \item Hence, $a_i^k$ is a
  %   best response to $\alpha(\varepsilon)$ for $\varepsilon$ small
  %   enough.
  \end{itemize}
\item   (3) implies (1):
  \begin{itemize}
  \item ``Cut off''--if necessary--the first few elements of the
    sequence of $\alpha(\varepsilon )$ such that for the remaining
    sequence $\alpha$ is a best reply to all $\alpha(\varepsilon
    )$. 
  \item $\eta_i^k(\varepsilon )=\alpha(\varepsilon )_i^k$ if $\alpha$ assigns zero probability to $a_i^k$
\item $\eta_i^k(\varepsilon )=min_l \alpha(\varepsilon )_i^l$ otherwise
% For actions $a_i^k$ that are played with zero probabiltiy
    % under $\alpha$, take $\eta_i^k=\alpha(\varepsilon )_i^k$. For
    % actions $a_i^k$ played with positive probability take a small
    % positive number, e.g. $\eta_i^k=min_l \alpha(\varepsilon
    % )_i^l$. 
  \item \dots % As $\varepsilon \ra 0$, so does the newly defined
  %   $\eta$.
  % \item note that actions that are not best replies are played only with probabiltiy $\eta_i^k$
  % \item $\alpha$ is a best response by each player in $G'(\eta)$ and therefore is an eqilibrium
  \end{itemize}
\end{itemize}
\qed

\framebreak

Some comments:
\begin{itemize}
\item the definition says: ``there is a sequence''; requiring that for \emph{any sequence}  of trembles converging to zero there is a sequence of equilibria leads to \emph{strict perfect equilibria}.  However, those do often not exist:
  \begin{ex}
     \begin{table}[h]
\centering
%\begin{tabular}{l|c|c|c}
 \begin{game}{2}{3}
    \> L   \>M  \>  R\\% \hline
T   \>1,1  \>1,0\>0,0  \\
B   \>1,1  \>0,0\>1,0
\end{game}
%\end{tabular}
\caption{no strict perfect equilibrium}
\label{tab:no_strict_perf_eq}
\end{table}
\begin{itemize}
\item $L$ is a strictly dominant strategy
\item  Check that $(T,L)$ and $(B,L)$ are both perfect equilibria but not strictly perfect 
  % \begin{itemize}
  % \item first take sequence where $\eta_2^M\great \eta_2^R$ $\Ra$ $T$ will be the unique best response in $G'(\eta)$ $\Ra$ $T$ will be played with certainty and $B$ with zero probability in limit
  %  \item   second take sequence where $\eta_2^M\less \eta_2^R$ $\Ra$ $B$ is the unique best response in $G'(\eta)$ $\Ra$ $B$ is played with probabiltiy 1 in limit
  % \end{itemize}
\end{itemize}
  \end{ex}

\item like most refinements perfection is sensitive to the formulation of the game:
  \begin{ex}\label{ex:perf_eq_sensitive}
        \begin{table}[h]
\centering
%\begin{tabular}{l|c|c|c}
 \begin{game}{3}{3}
      \> L   \>C   \>  R\\% \hline
T   \>1,1    \>0,0\>-1,-2  \\
M   \>0,0    \>0,0  \>0,-2\\
B   \> -2,-1 \>-2,0\>-2,-2
\end{game}
%\end{tabular}
\caption{sensitive to additional irrelevant actions}
\label{tab:sensitive_perf_eq}
\end{table}
\vspace*{-0.5cm}
\begin{itemize}
\item same game as first example plus strictly dominated actions B and R
\item Check that (M,C) is also perfect equilibrium in this game
 %  \begin{itemize}
%   \item 
% Take $\eta_t^R=1/t$ and $\eta_t^L=1/(2t)$ 
% \item M is unique b.r. in every $G'(\eta)$.
% \item also in limit $M$ is played with probability 1
% \item take $\eta_t^B=1/t$ and $\eta_t^T=1/(2t)$
% \item   M is unique b.r. in all $G'$
% \item also in limit M is played with probability 1
% \end{itemize}
\item (of course,  (T,L) remains a perfect equilibrium as well)
\end{itemize}
  \end{ex}
\end{itemize}

%   \begin{lemma}
%     In finite two player games, perfect equilibrium is the same as Nash equilibrium in not weakly dominated actions.
%   \end{lemma}
% Proof: skipped
% \textbf{Proof. }\\
% We use (without proof):\\
% \textbf{Lemma: }If a mixed strategy $\alpha_i'$ of player $i$ is not weakly dominated, then it is a best response to some \emph{completely mixed}  strategy $\alpha_{-i}'$.\qed
% \\
% \begin{itemize}
% \item let $\alpha_{-i}(\varepsilon)=(1-\varepsilon)\alpha_{-i}^*+\varepsilon
%   \alpha_{-i}'$
% \item use third definition of perfect equilibrium (see theorem before)
% \end{itemize}
%  % to see the first point: take a NE in not weakly dominated strategies $\alpha^*$ and now let $\alpha_{-i}(\varepsilon)=(1-\varepsilon)\alpha_{-i}^*+\varepsilon \alpha_{-i}'$ (where $\alpha'$ is as in the last bullet point) ; then $\alpha_i^*$ is a best response to $\alpha_{-i}(\varepsilon)$ which is completely mixed and converges to the NE as $\varepsilon\ra 0$. According to the third definition in \ref{thm:def_perf_eq} we get that $\alpha^*$ is a perfect equilibrium. (other direction is straightforward)
% \qed

\framebreak

\begin{itemize}
\item three player game: perfect equilibrium eliminates more than just equilibria in weekly dominated actions
 \begin{table}[h]
\centering
%\begin{tabular}{l|c|c|c}
 \begin{game}{2}{2}[][][A]
      \> L   \>R  \\% \hline
T   \>1,1,1    \>1,0,1 \\
D   \> 1,1,1 \>0,0,1
\end{game}
%\end{tabular}
\qquad
 \begin{game}{2}{2}[][][B]
      \> L   \>R  \\% \hline
T   \>1,0,0    \>0,1,2 \\
D   \> 0,1,0 \>1,0,0
\end{game}
\caption{a NE that is not perfect despite not being in weakly dominated strategies}
\label{tab:ne_not_perf_undom}
\end{table}
\begin{itemize}
\item (D,L,A) is a NE and none of the actions is weakly
  dominated
\item  not perfect: for small enough trembles of
  P2 and P3 T will always be better than D for P1; % \par
% (hidden reason: the profile where D is better than T (i.e. R and B) will only be
%   reached with probability $\varepsilon^2$ which is negligible for small $\varepsilon$)
\end{itemize}

\end{itemize}

\framebreak
A Nash equilibrium in completely mixed strategies is also a perfect equilibrium.\\Why? \begin{example}
\begin{table}
 \begin{game}{2}{2}
      \> H   \>T   \\% \hline
H   \>1,-1    \>-1,1  \\
T   \>-1,1    \>1,-1
\end{game}
\caption{Matching Pennies}
\end{table}
\end{example}

}

\frame{\ft{Review Questions}
  \begin{itemize}
  \item What is the idea behind refinements? (why might we need them?)
  \item What is the idea behind perfect equilibrium?
  \item How is perfect equilibrium defined? 
  \item What problems does perfect equilibrium have?
  \end{itemize}
reading: lecture notes by Hans Keiding p. 1-6; MSZ p. 262-267 (or OR 12.5.1)\\
*reading: lecture notes by Hans Keiding p. 7- 
}

\frame{\ft{Exercises}
  \begin{enumerate}
  \item     Show in the game below that (T,L) is not a perfect equilibrium and
    (B,R) is a perfect equilibrium.
    \begin{table}[h]
    \centering
      \begin{game}{3}{2}
        \>L\>R\\
        T \>2,1 \> 1,1\\
        M \>1,2 \> 0,2\\
        B \>0,0 \> 3,1
      \end{game}
    \end{table}
  \item Show that (D,L,A) is a Nash equilibrium but not a perfect
    equilibrium in the following 3 player game:
 \begin{table}[h]
\centering
%\begin{tabular}{l|c|c|c}
 \begin{game}{2}{2}[][][A]
      \> L   \>R  \\% \hline
        T   \>1,0,0    \>1,0,1 \\
        D \> 1,1,1 \>0,1,1
\end{game}
%\end{tabular}
\qquad
 \begin{game}{2}{2}[][][B]
      \> L   \>R  \\% \hline
        T   \>0,1,1    \>0,0,0 \\
        D \> 1,0,0 \>1,1,0
\end{game}
\end{table}
    % Player 2 strictly prefers L only if both other players tremble
    % but strictly prefers R if only P3 trembles. Hence, R is his
    % utility maximizing action if the trembles are small enough.
\item ${}^*$Show that no player uses a weakly dominated strategy in any perfect equilibrium.
  \end{enumerate}
}

% \section{proper equilibrium}

% \fr{\ft{Proper equilibrium}
% \begin{itemize}
% \item perfect equilibrum: players make some mistakes
% \item proper equilibrium: bigger mistakes are less likely
% \end{itemize}

% \begin{ex}
% \begin{table}[h]
% \centering
% %\begin{tabular}{l|c|c|c}
%  \begin{game}{3}{3}
%       \> L   \>C   \>  R\\% \hline
% T   \>1,1    \>0,0\>-1,-2  \\
% M   \>0,0    \>0,0  \>0,-2\\
% B   \> -2,-1 \>-2,0\>-2,-2
% \end{game}
% %\end{tabular}
% % \caption{sensitive to additional irrelevant actions}
% % \label{tab:sensitive_perf_eq}
% \end{table} 
% \begin{itemize}
% \item to show that (B,C) is perfect equilibrium, we used bigger trembles on R than on L
% \item R seems to be bigger mistake $\Ra$ reasonable?
% \end{itemize}
% \end{ex}


% \begin{definition}[$\varepsilon $-proper equilibrium] 
% $\alpha(\varepsilon )$ is an $\varepsilon
%   $-proper equilibrium for $\varepsilon \great 0$ if it is completely
%   mixed and
% $$u_i(\alpha(\varepsilon )_{-i},k)\less u_i(\alpha(\varepsilon )_{-i},l)\quad \Ra \quad \alpha(\varepsilon)_i^k \leq \varepsilon \alpha(\varepsilon )_i^l$$
% where $k, l\in A_i$.
% \end{definition}

% \begin{itemize}
% \item if $l$ is better than $k$, then the probability of $k$ being played is at most $\varepsilon $ times the probability of $l$ being played
% \item note: if $a_1$ is better than $a_2$ which in turn is better than $a_3$, then $a_2$ is at most $\varepsilon $ as likely as $a_1$ and $a_3$ is at most $\varepsilon^2$ as likely as $a_1$
% \end{itemize}


% % \begin{lemma}
% %   A $\varepsilon $-proper equilibrium exists for $0<\varepsilon <1$.
% % \end{lemma}
% % \textbf{Proof. }
% % notation:
% % \begin{itemize}
% % \item $m=max_i(\# A_i)$
% % \item strategy space we consider:
% % $$ S_i=\left\{\alpha_i\in\Delta A_i: \alpha_i(a_k)\geq \frac{\varepsilon^m}{m} \text{ for all }a_i \in A_i  \right\}$$
% % \item constrained response of player $i$
% %   \begin{multline*}
% %     r_i(\alpha_{-i})=\{\alpha_i\in S_i: \text{ if }u_i(\alpha_{-i},k)\less
% %     u_i(\alpha_{-i},l)\\
% %  \text{ then }\alpha_i(k)\leq \varepsilon \alpha_i(l)
% %     \text{ for all }k,l\in A_i \}
% %   \end{multline*}
% % \end{itemize}
% % We want to use Kakutani's fixed point theorem on $r=r_1\times \dots \times r_n$.
% % Convex- and compact-valuedness and closed graph property are shown as in proof of the Nash theorem.\\
% % To show: $r_i$ is non-empty for all $\alpha_{-i}$ \dots
% % . Nonemptiness is shown by constructing a member: Take $\rho(a_i,s)$ as the number of actions $b_i$ in $A_i$ such that $u_i(s_{-i},a_i)\less u_i(s_{-i},b_i)$. Then the following strategy belongs to $r_i$ 
%  % $$\hat{s}_i(a_i)= \frac{\varepsilon ^{\rho(a_i,s)}}{\sum_{a_i'\in A_i}\varepsilon^{\rho(a_i',s)}}\geq \frac{\varepsilon^m}{m}.$$
% %\qed

% \framebreak

% \begin{definition}[Proper equilibrium]
%   $\alpha$ is a \emph{proper equilibrium} if it is the limit of a
%   sequence $\alpha(\varepsilon )_{\varepsilon \ra0}$ of $\varepsilon
%   $-proper equilibrium.
% \end{definition}


% % \begin{theorem}
% % In every finite strategic form game a proper equilibrium exists.
% % \end{theorem}

% % \textbf{Proof. }skipped
% % Then we can take a sequence of $\varepsilon $-proper equilibria where $\varepsilon $ tends to zero. We can take a convergent subsequence of those equilibria. The limit of this subsequence is a proper equilibrium.
% %\qed

% \begin{ex}
% \begin{table}[h]
% \centering
% %\begin{tabular}{l|c|c|c}
%  \begin{game}{3}{3}
%       \> L   \>C   \>  R\\% \hline
% T   \>1,1    \> 0,0\>-1,-2  \\
% M   \>0,0    \>0,0  \>0,-2\\
% B   \> -2,-1 \>-2,0\>-2,-2
% \end{game}
% %\end{tabular}
% % \caption{sensitive to additional irrelevant actions}
% % \label{tab:sensitive_perf_eq}
% \end{table} 

%   \begin{itemize}
%   \item we showed that (M,C) is perfect eq
%   \item  (M,C) is not proper eq:
%     \begin{itemize}
%     \item tremble on R can only be $\varepsilon$ as big as the tremble on $L$% , since L is a
%       % better response than R.
%    \item but then T is a better response than M in any perturbed game%  and will therefore be played for sure in
%       % the limit. Hence,
%    \item   only (T,L) is a proper equilibrium
%     \end{itemize}

%   \end{itemize}
% \end{ex}

% \framebreak



% \begin{ex}
% \begin{itemize}
% \item  proper equilibrium can unfortunately also be sensitive to adding
%   strictly dominated actions
%    \begin{table}[h]
% \centering
% %\begin{tabular}{l|c|c}
% %\multicolumn{3}{c}{A}\\
%  \begin{game}{2}{2}[][][A]
%       \> L   \>R   \\ %\hline
% T   \>1,1,1    \> 0,0,1  \\
% D   \>0,0,1    \>0,0,1
% \end{game}
% %\end{tabular}
% \quad
% %\begin{tabular}{l|c|c}
% %\multicolumn{3}{c}{A}\\
%  \begin{game}{2}{2}[][][B]
%       \> L   \>  R\\ %\hline
% T   \>0,0,0    \> 0,0,0  \\
% D   \>0,0,0    \>  1,1,0  
% \end{game}
% %\end{tabular}
% % \caption{sensitive to adding strictly dominated actions}
% % \label{tab:sensitive_prop_eq}
% \end{table}
% \item B is strictly dominated for player 3
% \item in the game without action B only (T,L,x) is a proper equilibrium
% \item Show that  (D,R,A) is also a proper equilibrium in the game with B 
% % Just think of a
% %   tremble where each player plays his action from (B,R,A) with prob
% %   $1-\varepsilon $ and the other action with prob $\varepsilon $. Then
% %   player 1 and 2 are indifferent between their two actions for each
% %   $\varepsilon $. The limit is the equilibrium (B,R,A).
% \end{itemize}
% \end{ex}
% }



\end{document}